% GENERATED FILE. Edit canonical lessons, then run npm run generate:books.
% canonical-source-hash: fnv1a64:219c6c7cc755649f
% canonical-lessons: FR-C350-impact, FR-C350-repercussion, FR-C350-declencher, FR-C350-susciter, FR-C350-motif

\chapter{Les répercussions}
\label{ch:fr-les-repercussions}
\hlchaptermodality{350}

\begin{chapteropening}
\textbf{By the end of this chapter:} \emph{I can describe the impact of an event, its knock-on effects, what set it off, what it aroused and the motive behind it.}

Complete the last lesson of chapter 350: I can describe the impact of an event, its knock-on effects, what set it off, what it aroused and the motive behind it.
\end{chapteropening}

\section[l'impact]{l'impact — the impact}

\label{lesson:FR-C350-impact}



\begin{quote}
Before the new one: say the French for to influence, then the French for to generate.
\end{quote}

\subsection*{You'll want to know: l'impact}
\textbf{l'impact} — \textquotedblleft{}the impact\textquotedblright{}.

\begin{quote}
Ce projet a un impact sur tout le quartier.
\end{quote}

\begin{quote}
\emph{This project has an impact on the whole neighbourhood.}
\end{quote}

\emph{Avoir un impact sur} — to have an effect on. The final \emph{ct} is said. Masculine.

\begin{cousinweb}
\emph{impact} comes from Latin \emph{impāctus}, the participle of \emph{impingere} \textquotedblleft{}to dash against, to strike\textquotedblright{} (Lewis \& Short), built on \emph{pangere} \textquotedblleft{}to fix, to drive in\textquotedblright{}.
\end{cousinweb}

\begin{grammarlens}[title={What you've built}]
One more word for telling, explaining and describing at B1.
\end{grammarlens}

\subsection*{Your turn}
\begin{itemize}
\raggedright
  \item \emph{Say it:} \emph{l'impact}
  \item \emph{Say it:} \emph{l'impact}, once more
  \item \emph{Say it:} a sentence of your own with \emph{l'impact}
  \item \emph{Recall:} say the French for to influence, then the French for to generate, then say \emph{l'impact} again
\end{itemize}

\begin{tcolorbox}[breakable,colback=teal!4,colframe=teal!35!black,title={Before you move on}]
What does l'impact mean? (\textquotedblleft{}The impact\textquotedblright{}.) Say it once more.
\end{tcolorbox}
\section[la répercussion]{la répercussion — the repercussion, the knock-on effect}

\label{lesson:FR-C350-repercussion}



\begin{quote}
Before the new one: say the French for to generate, then the French for the impact.
\end{quote}

\subsection*{You'll want to know: la répercussion}
\textbf{la répercussion} — \textquotedblleft{}the repercussion, the knock-on effect\textquotedblright{}.

\begin{quote}
Cette décision aura des répercussions sur les prix.
\end{quote}

\begin{quote}
\emph{This decision will have repercussions on prices.}
\end{quote}

Mostly in the plural: \emph{des répercussions}.

\begin{cousinweb}
\emph{répercussion} is a learned borrowing from Latin \emph{repercussiō} \textquotedblleft{}a rebounding\textquotedblright{} (Lewis \& Short), from \emph{repercutere} \textquotedblleft{}to strike back\textquotedblright{}, built on \emph{quatere} \textquotedblleft{}to shake, to strike\textquotedblright{}.
\end{cousinweb}

\begin{grammarlens}[title={What you've built}]
One more word for telling, explaining and describing at B1.
\end{grammarlens}

\subsection*{Your turn}
\begin{itemize}
\raggedright
  \item \emph{Say it:} \emph{la répercussion}
  \item \emph{Say it:} \emph{la répercussion}, once more
  \item \emph{Say it:} a sentence of your own with \emph{la répercussion}
  \item \emph{Recall:} say the French for to generate, then the French for the impact, then say \emph{la répercussion} again
\end{itemize}

\begin{tcolorbox}[breakable,colback=teal!4,colframe=teal!35!black,title={Before you move on}]
What does la répercussion mean? (\textquotedblleft{}The repercussion, the knock-on effect\textquotedblright{}.) Say it once more.
\end{tcolorbox}
\section[déclencher]{déclencher — to set off, to trigger}

\label{lesson:FR-C350-declencher}



\begin{quote}
Before the new one: say the French for the impact, then the French for the repercussion.
\end{quote}

\subsection*{You'll want to know: déclencher}
\textbf{déclencher} — \textquotedblleft{}to set off, to trigger\textquotedblright{}.

\begin{quote}
Une petite remarque a déclenché la colère de mon père.
\end{quote}

\begin{quote}
\emph{A small remark set off my father's anger.}
\end{quote}

Of alarms too: \emph{l'alarme s'est déclenchée}.

\begin{cousinweb}
\emph{déclencher} is built from \emph{dé-} and \emph{clenche} \textquotedblleft{}latch\textquotedblright{}, which comes from a Frankish \emph{klinka} \textquotedblleft{}latch\textquotedblright{} (EtymDB): you lift the latch and things start moving.
\end{cousinweb}

\begin{grammarlens}[title={What you've built}]
One more word for telling, explaining and describing at B1.
\end{grammarlens}

\subsection*{Your turn}
\begin{itemize}
\raggedright
  \item \emph{Say it:} \emph{déclencher}
  \item \emph{Say it:} \emph{déclencher}, once more
  \item \emph{Say it:} a sentence of your own with \emph{déclencher}
  \item \emph{Recall:} say the French for the impact, then the French for the repercussion, then say \emph{déclencher} again
\end{itemize}

\begin{tcolorbox}[breakable,colback=teal!4,colframe=teal!35!black,title={Before you move on}]
What does déclencher mean? (\textquotedblleft{}To set off, to trigger\textquotedblright{}.) Say it once more.
\end{tcolorbox}
\section[susciter]{susciter — to arouse, to give rise to}

\label{lesson:FR-C350-susciter}



\begin{quote}
Before the new one: say the French for the repercussion, then the French for to set off.
\end{quote}

\subsection*{You'll want to know: susciter}
\textbf{susciter} — \textquotedblleft{}to arouse, to give rise to\textquotedblright{}.

\begin{quote}
Sa décision a suscité beaucoup de questions.
\end{quote}

\begin{quote}
\emph{Her decision raised a lot of questions.}
\end{quote}

Formal: \emph{susciter la curiosité} — to arouse curiosity.

\begin{cousinweb}
\emph{susciter} is a learned borrowing from Latin \emph{suscitāre} \textquotedblleft{}to lift up, to rouse\textquotedblright{} (EtymDB; Lewis \& Short), built on \emph{citāre} \textquotedblleft{}to set in motion\textquotedblright{}.
\end{cousinweb}

\begin{grammarlens}[title={What you've built}]
One more word for telling, explaining and describing at B1.
\end{grammarlens}

\subsection*{Your turn}
\begin{itemize}
\raggedright
  \item \emph{Say it:} \emph{susciter}
  \item \emph{Say it:} \emph{susciter}, once more
  \item \emph{Say it:} a sentence of your own with \emph{susciter}
  \item \emph{Recall:} say the French for the repercussion, then the French for to set off, then say \emph{susciter} again
\end{itemize}

\begin{tcolorbox}[breakable,colback=teal!4,colframe=teal!35!black,title={Before you move on}]
What does susciter mean? (\textquotedblleft{}To arouse, to give rise to\textquotedblright{}.) Say it once more.
\end{tcolorbox}
\section[le motif]{le motif — the motive, the reason; the pattern}

\label{lesson:FR-C350-motif}



\begin{quote}
Before the new one: say the French for to set off, then the French for to arouse.
\end{quote}

\subsection*{You'll want to know: le motif}
\textbf{le motif} — \textquotedblleft{}the motive, the reason; the pattern\textquotedblright{}.

\begin{quote}
Quel est le motif de votre visite ?
\end{quote}

\begin{quote}
\emph{What is the reason for your visit?}
\end{quote}

Also a pattern on cloth: \emph{un motif à fleurs}.

\begin{cousinweb}
\emph{motif} comes from Late Latin \emph{mōtīvus} \textquotedblleft{}serving to move\textquotedblright{} (EtymDB), from \emph{movēre} \textquotedblleft{}to move\textquotedblright{} (Lewis \& Short): what moves you to act.
\end{cousinweb}

\begin{grammarlens}[title={What you've built}]
One more word for telling, explaining and describing at B1.
\end{grammarlens}

\subsection*{Your turn}
\begin{itemize}
\raggedright
  \item \emph{Say it:} \emph{le motif}
  \item \emph{Say it:} \emph{le motif}, once more
  \item \emph{Say it:} a sentence of your own with \emph{le motif}
  \item \emph{Recall:} say the French for to set off, then the French for to arouse, then say \emph{le motif} again
\end{itemize}

\begin{tcolorbox}[breakable,colback=teal!4,colframe=teal!35!black,title={Before you move on}]
What does le motif mean? (\textquotedblleft{}The motive, the reason; the pattern\textquotedblright{}.) Say it once more.
\end{tcolorbox}
